% ==============================================================
\section{Proposed Method: Multi-Dimensional Apriori-Window Algorithm}
\label{sec:method}
% ==============================================================

We propose the \textbf{Multi-Dimensional Apriori-Window algorithm}
to solve the problem in \S\ref{sec:dense_detection}.

% ==============================================================
\subsection{1D Multiset Reduction via Fixed Spatial Start Points}\label{sec:binsearch}
% ==============================================================

For a pattern $P$ and a fixed spatial start point
$\bm{s}_{\mathrm{sp}}$,
we construct the temporal multiset
$L_P(\bm{s}_{\mathrm{sp}})$
from Definition~\ref{def:collapsed_multiset_nd}.
This is obtained by distributing each occurrence transaction of $P$
to all spatial start points whose window contains that transaction.

\paragraph{Distribution to spatial start points.}
When a transaction containing pattern $P$
resides at discrete coordinates
\[
  \bm{b} = (b^{(0)}, b^{(1)}, \ldots, b^{(d-1)}),
\]
the $k$-th component ($k \geq 1$)
of a spatial start point containing this transaction satisfies
\[
  s_k \in [\max(0,\, b^{(k)} - W_k),\; b^{(k)}].
\]
Therefore, one occurrence contributes to at most
\[
  W_{\mathrm{sp}} = \prod_{k=1}^{d-1} (W_k + 1)
\]
spatial start points.

\paragraph{Construction cost.}
Let $N_P$ denote the number of occurrence transactions of pattern $P$.
Construction costs $O(N_P \cdot W_{\mathrm{sp}})$;
including sorting, the total is $O(\widetilde N_P \log \widetilde N_P)$
where $\widetilde N_P = N_P \cdot W_{\mathrm{sp}}$.

% ==============================================================
\subsection{Duplicate-Aware Interval Solver on the $t$-Axis}\label{sec:t_stride}
% ==============================================================

For a fixed spatial start point $\bm{s}_{\mathrm{sp}}$,
we compute dense intervals from the sorted multiset
$L = L_P(\bm{s}_{\mathrm{sp}})$.
In two or more dimensions,
multiple transactions may share the same timestamp,
producing duplicates in $L$.
Consequently, the simple ``$+1$ jump'' of the one-dimensional case
cannot be directly applied.

The key idea is a two-branch stride
governed by the comparison of
the window count $c$ and the threshold $\theta$
(Figure~\ref{fig:t_stride}).
When $c < \theta$ (non-dense),
the left endpoint $l$ jumps to the position
where the next event enters the window,
skipping unnecessary scans
(Eq.~\eqref{eq:jump_empty}, Figure~\ref{fig:t_stride}(b)).
When $c \geq \theta$ (dense),
the surplus $c - \theta$ determines
a contiguous block $[l, l_{\mathrm{dense}}]$ of left endpoints
that are all guaranteed to be dense,
and this entire block is output as a dense interval in bulk
(Eq.~\eqref{eq:jump_dense}, Figure~\ref{fig:t_stride}(a)).
The correctness of this bulk processing
is guaranteed by Lemma~\ref{lem:stride}.

Specifically, the left endpoint update in the non-dense case is
\begin{equation}\label{eq:jump_empty}
  l \leftarrow \max(l+1,\; L[\mathrm{next}] - W_0)
\end{equation}
where $L[\mathrm{next}]$ is the first event
outside the right boundary of the current window.
In the dense case, the block right boundary is
\begin{equation}\label{eq:jump_dense}
  l_{\mathrm{dense}} = L[l_{\mathrm{lo}} + (c - \theta)]
\end{equation}
where $l_{\mathrm{lo}}$ is the first position in $L$
that is $\geq l$.
The dense interval's right endpoint is extended to
$l_{\mathrm{dense}} + W_0$,
and the left endpoint is updated to
$l \leftarrow l_{\mathrm{dense}} + 1$.

\begin{lemma}[Correctness of the dense block]\label{lem:stride}
For $l_{\mathrm{dense}}$ defined by Eq.~\eqref{eq:jump_dense},
for any $l' \in [l,\, l_{\mathrm{dense}}]$,
\[
  |\{t \in L : l' \leq t \leq l' + W_0\}| \geq \theta
\]
holds.
\end{lemma}

The proof is given in Appendix~\ref{sec:proofs}.

\begin{figure}[t]
\centering
\includegraphics[width=0.85\columnwidth]{fig/fig3_t_stride.pdf}
\caption{Duplicate-aware $t$-axis stride.
Given a sorted multiset
$L = \{\!\{1, 2, 2, 3, 3, 8\}\!\}$
with $W_0{=}3$ and $\theta{=}4$.
(a)~When $c \geq \theta$,
the surplus determines how far the left endpoint
can stride in bulk (Eq.~\eqref{eq:jump_dense}).
(b)~When $c < \theta$,
the window jumps to the next position
where a new event enters (Eq.~\eqref{eq:jump_empty}).
Window color gradient (dark$\to$light blue)
indicates the progression of the left endpoint $l$.}
\label{fig:t_stride}
\end{figure}

\begin{algorithm}[t]
\caption{Duplicate-aware dense interval computation for a fixed spatial start point}
\label{alg:t_stride}
\begin{algorithmic}[1]
\Require Sorted multiset $L$, temporal window width $W_0$, threshold $\theta$
\Ensure Dense interval sequence $\mathcal{R}$
\State $\mathcal{R} \leftarrow \emptyset$
\State $l \leftarrow \max(0,\, L[\theta-1] - W_0)$
\State $\mathit{open} \leftarrow \mathrm{false}$
\While{$l \leq L[\mathrm{last}]$}
  \State $l_{\mathrm{lo}} \leftarrow \mathrm{lower\_bound}(L, l)$
  \State $l_{\mathrm{hi}} \leftarrow \mathrm{upper\_bound}(L, l + W_0)$
  \State $c \leftarrow l_{\mathrm{hi}} - l_{\mathrm{lo}}$
  \If{$c < \theta$}
    \If{$\mathit{open}$}
      \State Append current interval to $\mathcal{R}$
      \State $\mathit{open} \leftarrow \mathrm{false}$
    \EndIf
    \If{$l_{\mathrm{hi}} = |L|$} \textbf{break} \EndIf
    \State $l \leftarrow \max(l+1,\; L[l_{\mathrm{hi}}] - W_0)$
           \Comment{Eq.~\eqref{eq:jump_empty}}
  \Else
    \State $l_{\mathrm{dense}} \leftarrow L[l_{\mathrm{lo}} + (c-\theta)]$
           \Comment{Eq.~\eqref{eq:jump_dense}}
    \State $r_{\mathrm{block}} \leftarrow l_{\mathrm{dense}} + W_0$
    \State Start or extend current interval with $(l, r_{\mathrm{block}})$
    \State $\mathit{open} \leftarrow \mathrm{true}$
    \State $l \leftarrow l_{\mathrm{dense}} + 1$
  \EndIf
\EndWhile
\If{$\mathit{open}$}
  \State Append current interval to $\mathcal{R}$
\EndIf
\State \Return $\mathcal{R}$
\end{algorithmic}
\end{algorithm}

% ==============================================================
\subsection{Candidate Interval Restriction via Apriori Pruning ($k \geq 2$)}\label{sec:alg_k2}
% ==============================================================

Applying Lemma~\ref{lem:apriori} recursively,
any dense-start component of a size-$k$ pattern $P$
is contained within the dense intervals
of each $(k{-}1)$-subset of $P$.
Using this property,
for candidate patterns $P$ we use
\emph{intersections of $(k{-}1)$-subset interval maps}.

\paragraph{Candidate interval map.}
Let the dense region of pattern $Q$ be
\[
  \mathcal{R}_\theta(Q)
  = \{(\bm{s}_{\mathrm{sp}}, [a,e])\}.
\]
Write $\mathcal{R}_\theta(Q)[\bm{s}_{\mathrm{sp}}]$
for the temporal interval sequence
belonging to the same spatial start point $\bm{s}_{\mathrm{sp}}$.
The candidate interval map for a size-$k$ pattern $P$ is then defined as
\begin{equation}\label{eq:candidate_cells}
  \mathcal{C}(P,\bm{s}_{\mathrm{sp}})
  =
  \bigcap_{Q \in \binom{P}{k-1}}
  \mathcal{R}_\theta(Q)[\bm{s}_{\mathrm{sp}}].
\end{equation}

Here the intersection is the standard interval intersection
on temporal intervals, and $\binom{P}{k-1}$ denotes the set of all
$(k{-}1)$-subsets of $P$.
If $\mathcal{C}(P,\bm{s}_{\mathrm{sp}})=\emptyset$
for all spatial start points,
$P$ cannot have any dense interval and is pruned.
For $k=2$, $\binom{P}{1}$ reduces to singletons,
so this is a generalization of the singleton intersection.

\paragraph{Exact re-evaluation.}
Only when $\mathcal{C}(P,\bm{s}_{\mathrm{sp}}) \neq \emptyset$,
we scan the transaction sequence once
to construct $L_P(\bm{s}_{\mathrm{sp}})$
and run the candidate-interval version of
Algorithm~\ref{alg:t_stride}
over $\mathcal{C}(P,\bm{s}_{\mathrm{sp}})$.

% ==============================================================
\subsection{Overall Algorithm and Complexity}\label{sec:complexity}
% ==============================================================

\begin{algorithm}[t]
\caption{Multi-Dimensional Apriori-Window (interval-based implementation)}
\label{alg:md_apriori}
\begin{algorithmic}[1]
\Require $d$-dimensional transaction sequence $\mathcal{D}$, $(\bm{W}, \theta)$
\Ensure Dense region $\mathcal{R}_\theta(P)$ for each pattern $P$
\ForAll{$i \in \mathcal{I}$}
  \State $L_i(\bm{s}_{\mathrm{sp}}) \leftarrow$ Construct temporal multiset per spatial start point
  \State $\mathcal{R}_\theta(\{i\}) \leftarrow$ Run Algorithm~\ref{alg:t_stride} for each $\bm{s}_{\mathrm{sp}}$
\EndFor
\State $\mathcal{F}_1 \leftarrow \{\{i\} : \mathcal{R}_\theta(\{i\}) \neq \emptyset\}$
\State $k \leftarrow 2$
\While{$\mathcal{F}_{k-1} \neq \emptyset$}
  \State \texttt{cands} $\leftarrow$ \Call{GenerateCandidates}{$\mathcal{F}_{k-1}$}
  \State \texttt{cands} $\leftarrow$ \Call{PruneCandidates}{\texttt{cands}, $\mathcal{F}_{k-1}$}
  \State $\mathcal{F}_k \leftarrow \emptyset$
  \ForAll{$P \in$ \texttt{cands}}
    \State $\mathcal{C}(P,\bm{s}_{\mathrm{sp}}) \leftarrow \bigcap_{Q \in \binom{P}{k-1}}\mathcal{R}_\theta(Q)[\bm{s}_{\mathrm{sp}}]$
           \Comment{Eq.~\eqref{eq:candidate_cells}}
    \If{$\mathcal{C}(P,\bm{s}_{\mathrm{sp}})=\emptyset$ for all $\bm{s}_{\mathrm{sp}}$}
      \State \textbf{continue}
    \EndIf
    \State $L_P(\bm{s}_{\mathrm{sp}}) \leftarrow$ Construct temporal multiset for pattern $P$
    \State $\mathcal{R}_\theta(P) \leftarrow$ Run candidate-interval duplicate-aware solver for each $\bm{s}_{\mathrm{sp}}$
    \If{$\mathcal{R}_\theta(P) \neq \emptyset$}
      \State $\mathcal{F}_k \leftarrow \mathcal{F}_k \cup \{P\}$
    \EndIf
  \EndFor
  \State $k \leftarrow k + 1$
\EndWhile
\State \Return $\{\mathcal{R}_\theta(P)\}$
\end{algorithmic}
\end{algorithm}

\begin{table}[t]
\centering
\caption{Computational complexity of the proposed method.
$N_P$: number of transactions containing pattern $P$,
$W_{\mathrm{sp}} = \prod_{k=1}^{d-1} (W_k + 1)$,
$\widetilde N_P = N_P W_{\mathrm{sp}}$,
$I_P$: total number of intervals in the candidate interval map.}
\label{tab:complexity}
\small
\begin{tabular}{lll}
\toprule
Step & Complexity \\
\midrule
Construction of $L_P$ & $O(N_P \cdot W_{\mathrm{sp}})$ \\
Sorting of $L_P$ & $O(\widetilde N_P \log \widetilde N_P)$ \\
Temporal interval & \multirow{2}{*}{$O\!\left(\sum_{\bm{s}_{\mathrm{sp}}}|L_P(\bm{s}_{\mathrm{sp}})| \log |L_P(\bm{s}_{\mathrm{sp}})|\right)$} \\
\quad computation & \\
Candidate interval & \multirow{2}{*}{$O(I_P)$} \\
\quad map intersection & \\
\bottomrule
\end{tabular}
\end{table}

Therefore, the overall complexity
faithful to the current implementation is
\[
  O\!\left(
    \sum_{P \in \mathcal{F}}
      \bigl(
        N_P W_{\mathrm{sp}}
        +
        \widetilde N_P \log \widetilde N_P
        +
        I_P
      \bigr)
  \right)
\]
where $\mathcal{F}$ is the set of candidate patterns
actually evaluated.
When applied to real data,
discretization of continuous coordinates is required
as a preprocessing step (\S\ref{sec:experiments}).
